\section{个人思考}

\subsection{最重要的理解}

Vision-Zero 最值得关注的地方不是“用强化学习训练 VLM”，而是它把训练数据生成、难度控制和奖励计算统一成了一个游戏环境。很多 VLM 后训练工作把核心精力放在构造更好的题库、更长的 CoT 或更精细的人工验证器上；Vision-Zero 则把问题往前推了一步：先设计一个能自动产生有用监督信号的环境，再让模型在环境中通过交互生成数据。

从这个角度看，Vision-Zero 与 MM-Zero 有相似之处：二者都希望减少人工标注依赖，并让模型自己产生训练信号。但二者路径不同。MM-Zero 更强调 Proposer--Coder--Solver 闭环，通过代码渲染构造可验证视觉任务；Vision-Zero 更强调多智能体不完全信息博弈，通过角色冲突和投票结果构造奖励。前者更像“自动造题”，后者更像“自动比赛”。

\subsection{方法优点}

第一个优点是任务能力重合度高。普通玩家给线索需要识别图像中的关键元素，卧底伪装需要从语言线索中反推视觉内容，投票阶段需要对比图像和文本描述。这些操作与 VLM 的真实应用任务高度相关，不是与下游任务相距很远的抽象游戏。

第二个优点是奖励设计相对自然。卧底收到票数、普通玩家是否投对，都能由游戏规则直接给出，不需要人工评价线索“好不好”。这使 RLVR 的可验证性来自交互结果，而不是人工写死的答案。

第三个优点是交替训练机制比较务实。纯自博弈常见的问题是双方一起变强或一起变差，奖励归因不稳定；纯 RLVR 的问题是固定任务很快被学会。Iterative-SPO 用决策准确率和 \texttt{N/A} 比例判断当前阶段是否饱和，动态切换训练对象，工程上比完全开放式自博弈更容易控制。

\subsection{潜在局限}

Vision-Zero 说自己支持任意图像，但实际训练仍需要构造“普通玩家看到的图像”和“卧底看到的图像”之间的非对称信息。CLEVR 可以用渲染器生成，图表可以解析成 JSON 后重绘，真实图像可以用图像编辑模型修改；这些都不是完全免费的前提。对专业图像领域来说，可靠编辑和差异控制可能成为新瓶颈。

另一个风险是游戏目标与真实任务之间仍然存在偏差。模型为了赢游戏，可能学到的是“如何说不容易被怀疑的线索”和“如何利用其他玩家表达习惯”，而不一定完全等价于更强的视觉理解。论文用多任务评测证明了正迁移，但如果长期训练，仍需要警惕策略性语言模式压过真实视觉能力。

还有一个问题是多智能体训练的可解释性。胜负奖励看起来可验证，但某个玩家输赢可能由很多因素共同造成：图像难度、前面线索质量、卧底伪装水平、普通玩家投票能力、模型采样随机性等。RAE 和 group normalization 缓解了部分问题，但不能完全解决信用分配。

\subsection{对自进化研究的启发}

这篇论文说明，自进化系统的关键不只是“让模型自己生成数据”，而是要让生成过程具有以下性质：

\begin{enumerate}
    \item 生成的数据必须和目标能力有真实重合，而不是只是在形式上像训练任务。
    \item 奖励必须可自动计算，并且尽量不被简单捷径利用。
    \item 难度必须随模型能力提升而提升，否则训练会很快饱和。
    \item 训练分布必须保持多样，否则模型会收敛到少数容易赢的局部策略。
\end{enumerate}

Vision-Zero 的“普通玩家--卧底--投票者”结构，把数据难度和模型能力绑定起来。模型越会伪装，投票越难；模型越会识别，线索阶段就必须产生更高难度对抗。这种互相追赶的结构，是比静态题库更有潜力的自进化机制。

\subsection{复现建议}

如果要复现或扩展 Vision-Zero，我会优先做一个最小闭环，而不是一开始就复现全量多域实验：

\begin{enumerate}
    \item 先选 CLEVR 或其他可程序化渲染环境，保证图像差异和真值完全可控。
    \item 固定 3--4 个普通玩家和 1 个卧底，先实现 clue stage、decision stage、投票解析和奖励计算。
    \item 单独记录卧底胜率、普通玩家投票准确率、\texttt{N/A} 比例、线索长度和无效格式率，避免只看下游 benchmark。
    \item 先做 RAE 的有无消融，因为论文显示不带 RAE 会造成明显负迁移。
    \item 再考虑把图像类型扩展到图表或真实图像，并检查编辑失败、答案泄露和线索模板化问题。
\end{enumerate}

\subsection{后续问题}

\begin{itemize}
    \item Vision-Zero 长期迭代后会继续提升，还是最终收敛到某种语言策略均衡？
    \item 卧底看到空白图像和看到修改图像，这两种设置对能力迁移的影响是否不同？
    \item 如果不同角色使用不同模型，例如强模型当裁判、弱模型当玩家，训练稳定性会不会更好？
    \item 游戏奖励能否加入更细的视觉一致性验证，减少纯语言伪装策略？
    \item 这个框架能否扩展到视频、多图检索或 3D 交互环境，而不仅是单图差异？
\end{itemize}
